% =====================================================================
%  1. Introducción
% =====================================================================
\section{Introducción}

\subsection{Presentación del documento}

\begin{instruccion}
Explique el objetivo del plan de prueba, su estructura y, de forma general, el contenido de
cada apartado. Indique además a quién va dirigido el documento y en qué momento del proyecto
se emite. Conviene señalar explícitamente qué partes de \norma\ se han tomado como referencia
y cuáles no, para no dar a entender una conformidad que no se ha demostrado.
\end{instruccion}

\campo{Redacte aquí la presentación del documento}

\begin{notanormativa}[Relación entre las partes de la norma y este documento]
\normap{2} describe los \emph{procesos} de prueba; \normap{3}, la \emph{documentación} que
esos procesos producen; \normap{4}, las \emph{técnicas} de diseño de pruebas. Por lo tanto,
el contenido de este plan \textbf{no sustituye} a los registros que se generan al diseñar,
ejecutar, controlar y cerrar las pruebas: los complementa.\sidiag{ La
Figura~\ref{fig:estructura} muestra esa relación y señala qué apartados de este documento se
apoyan en cada parte.}
\end{notanormativa}

\ifmostrardiagramas
% =====================================================================
%  diagrams/estructura-29119.tex
%  Figura: estructura general de ISO/IEC/IEEE 29119 y su relación
%          con las partes de este documento.
%
%  CORRECCIONES respecto del diagrama original («normal general.png»):
%   1. El original colgaba «Procesos organizacionales · política y
%      estrategia» de la Parte 1. Los procesos —incluido el
%      organizacional— pertenecen a la Parte 2; la Parte 1 aporta
%      conceptos generales y vocabulario.
%   2. Se elimina toda sugerencia de que una parte sustituye o contiene
%      a otra: son partes complementarias de una misma norma.
%   3. Se añade la correspondencia explícita con las secciones y anexos
%      de este documento, que el original no representaba.
%
%  Nota de mantenimiento: los nodos se colocan con «positioning»
%  (below=… of …) en lugar de coordenadas absolutas, de modo que al
%  cambiar el texto de un cuadro los demás se desplacen y nunca se
%  superpongan.
% =====================================================================
\begin{figure}[htbp]
\centering
\begin{tikzpicture}[
  font=\scriptsize,
  node distance=4mm,
  raiz/.style   ={draw=uvazul, very thick, fill=uvazul!12, rounded corners=3pt,
                  text width=5.0cm, align=center, inner sep=4pt},
  parte/.style  ={draw=uvazul, thick, fill=uvazul!7, rounded corners=2pt,
                  text width=2.35cm, align=center, inner sep=4pt},
  sub/.style    ={draw=gray!70, fill=gray!8, rounded corners=2pt,
                  text width=2.35cm, align=center, inner sep=4pt},
  doc/.style    ={draw=instrverde, thick, fill=instrverde!8, rounded corners=2pt,
                  text width=3.7cm, align=center, inner sep=4pt},
  fl/.style     ={-{Latex[length=4pt]}, draw=uvazul!80},
  mapa/.style   ={-{Latex[length=4pt]}, draw=instrverde!85, dashed}
]

% ------------------------------- raíz --------------------------------
\node[raiz] (n) {\textbf{ISO/IEC/IEEE 29119}\\ Prueba de software (norma multiparte)};

% ------------------------------ partes -------------------------------
\node[parte, below=10mm of n] (p3) {\textbf{Parte 3}\\ Documentación de prueba};
\node[parte, left=6mm of p3]  (p2) {\textbf{Parte 2}\\ Procesos de prueba};
\node[parte, left=6mm of p2]  (p1) {\textbf{Parte 1}\\ Conceptos generales y vocabulario};
\node[parte, right=6mm of p3] (p4) {\textbf{Parte 4}\\ Técnicas de prueba};
\node[parte, right=6mm of p4] (p5) {\textbf{Parte 5}\\ Pruebas por palabras clave};

\foreach \d in {p1,p2,p3,p4,p5} { \draw[fl] (n) -- (\d); }

% --------------------------- contenido -------------------------------
% Parte 1
\node[sub, below=of p1] (v1) {Vocabulario y conceptos que usan las demás partes};
\draw[fl] (p1) -- (v1);

% Parte 2: las tres capas de procesos
\node[sub, below=of p2] (c1) {Capa organizacional};
\node[sub, below=of c1] (c2) {Gestión de la prueba};
\node[sub, below=of c2] (c3) {Procesos dinámicos};
\draw[fl] (p2) -- (c1); \draw[fl] (c1) -- (c2); \draw[fl] (c2) -- (c3);

% Parte 3: tipos de documentación
\node[sub, below=of p3] (d1) {Documentos de gestión};
\node[sub, below=of d1] (d2) {Especificaciones de nivel};
\node[sub, below=of d2] (d3) {Registros de ejecución e incidentes};
\draw[fl] (p3) -- (d1); \draw[fl] (d1) -- (d2); \draw[fl] (d2) -- (d3);

% Parte 4: familias de técnicas
\node[sub, below=of p4] (t1) {Basadas en especificación};
\node[sub, below=of t1] (t2) {Basadas en estructura};
\node[sub, below=of t2] (t3) {Basadas en experiencia};
\draw[fl] (p4) -- (t1); \draw[fl] (t1) -- (t2); \draw[fl] (t2) -- (t3);

% Parte 5
\node[sub, below=of p5] (k1) {Enfoque opcional; se usa sólo si el proyecto lo adopta};
\draw[fl] (p5) -- (k1);

% ------------------- correspondencia con este documento --------------
\node[fit=(v1)(c3)(d3)(t3)(k1), inner sep=0pt, draw=none] (todos) {};

\node[align=center, font=\scriptsize\bfseries, text=instrverde,
      below=7mm of todos.south] (rot) {ESTE DOCUMENTO};

\node[doc, below=3mm of rot] (e2) {\textbf{Parte II}\\ Reporte de finalización};
\node[doc, left=5mm of e2]   (e1) {\textbf{Parte I}\\ Plan de prueba (secciones 1--9)};
\node[doc, right=5mm of e2]  (e3) {\textbf{Anexos A--G}\\ casos, procedimientos, ejecuciones, incidentes, informes, entorno, trazabilidad};

\draw[mapa] (c2.west) to[out=190,in=100] (e1.north);
\draw[mapa] (d1.east) to[out=0,in=90]    (e2.north east);
\draw[mapa] (d3.east) to[out=350,in=110] (e3.north west);
\draw[mapa] (t3.south) to[out=280,in=80] (e3.north);

% ------------------------------ leyenda ------------------------------
\node[align=left, font=\scriptsize, text width=15.2cm, below=5mm of e2] (leg)
  {\textbf{Cómo leer la figura.} Las flechas azules indican \emph{qué contiene} cada parte de la
   norma; las verdes discontinuas, \emph{dónde} se refleja ese contenido en este documento.
   Ninguna parte sustituye a otra: el plan no reemplaza a los registros que producen los procesos
   dinámicos, ni las técnicas de la Parte~4 son obligatorias en su totalidad.};

\end{tikzpicture}
\caption[Estructura de ISO/IEC/IEEE 29119 y su relación con este documento]%
{Estructura general de \norma\ y correspondencia de sus partes con las secciones y anexos de
este documento.}
\label{fig:estructura}
\end{figure}

\begin{observacion}[Diagrama de estructura de la norma]
\obsproblema{El diagrama de referencia colgaba «Procesos organizacionales · política y
estrategia» de la Parte~1, y sugería una jerarquía en la que unas partes sustituían a otras.}
\obscambio{Los procesos —incluido el organizacional— cuelgan de la Parte~2; la Parte~1 aporta
conceptos y vocabulario. Se añade la correspondencia con las secciones y anexos de este
documento y una leyenda que aclara que ninguna parte reemplaza a otra.}
\obsfuente{OBS §1 · DIAG}
\end{observacion}


La Figura~\ref{fig:estructura} presenta la relación conceptual entre las partes de \norma\ y
las secciones y anexos de esta plantilla. Debe leerse como un mapa de correspondencias, no
como una jerarquía de sustitución: ninguna parte de la norma reemplaza a otra, y disponer de
este plan no acredita por sí solo la aplicación de los procesos que la norma describe.

% =====================================================================
%  diagrams/procesos-prueba.tex
%  Figura: flujo del proceso de prueba en tres capas, con sus bucles
%          de retroalimentación.
%
%  CORRECCIONES respecto del diagrama original («Procesos de prueba.png»):
%   1. El original rotulaba «verificar correcciones» como prueba de
%      regresión. Aquí se representan por separado la prueba de
%      CONFIRMACIÓN (el defecto quedó corregido) y la de REGRESIÓN (el
%      cambio no dañó lo que ya funcionaba).
%   2. El original se leía como una secuencia lineal. Aquí el monitoreo
%      y control aparece como proceso CONCURRENTE y se dibujan los tres
%      bucles de retroalimentación: control, replanificación y mejora.
%   3. El original situaba el entorno y los datos dentro de la
%      preparación de pruebas; en la edición 2021 constituyen un proceso
%      propio (ES1/ES2), y así se representa.
%   4. El original sólo contemplaba el análisis de riesgos del proyecto;
%      la planificación abarca riesgos de producto y de proyecto.
% =====================================================================
\begin{figure}[htbp]
\centering
\begin{tikzpicture}[
  font=\scriptsize,
  proc/.style   ={draw=uvazul, thick, fill=uvazul!8, rounded corners=2pt,
                  text width=2.8cm, align=center, minimum height=1.0cm, inner sep=3pt},
  procm/.style  ={draw=polinaranja, very thick, fill=polfondo, rounded corners=2pt,
                  text width=2.9cm, align=center, minimum height=1.5cm, inner sep=3pt},
  salida/.style ={draw=gray!70, fill=gray!8, rounded corners=2pt,
                  text width=2.6cm, align=center, minimum height=0.7cm, inner sep=3pt},
  conf/.style   ={draw=instrverde, thick, fill=instrverde!8, rounded corners=2pt,
                  text width=2.7cm, align=center, minimum height=0.9cm, inner sep=3pt},
  fl/.style     ={-{Latex[length=4pt]}, draw=uvazul!85, thick},
  fb/.style     ={-{Latex[length=4pt]}, draw=polinaranja, thick, dashed},
  mej/.style    ={-{Latex[length=4pt]}, draw=instrverde, thick, dashdotted},
  et/.style     ={font=\tiny, inner sep=1pt, fill=white}
]

% ============================ CAPA ORGANIZACIONAL =====================
\node[proc] (ot) at (-4.6,0) {\textbf{Proceso organizacional}\\ OT1 desarrollar · OT2 monitorear · OT3 actualizar};
\node[salida] (pol) at (-0.6,0) {Política de prueba\\ y estrategia organizacional};
\draw[fl] (ot) -- (pol);

% ============================ CAPA DE GESTIÓN =========================
\node[proc] (tp) at (-5.6,-2.6) {\textbf{Planificación}\\ TP1--TP9\\ \emph{riesgos de producto y de proyecto}};
\node[salida] (plan) at (-5.6,-4.3) {Plan de prueba};
\draw[fl] (tp) -- (plan);

\node[procm] (tmc) at (-0.4,-3.3) {\textbf{Monitoreo y control}\\ TMC1--TMC4\\ \emph{concurrente con la ejecución}};

\node[proc] (tc) at (4.6,-3.3) {\textbf{Finalización}\\ TC1 archivar · TC2 limpiar entorno · TC3 lecciones · TC4 reportar};

\draw[fl] (pol) to[out=250,in=70] node[et,pos=0.5]{gobierna} (tp.north);

% ============================ CAPA DINÁMICA ===========================
\node[proc] (es) at (-6.0,-6.6) {\textbf{Entorno y datos}\\ ES1 establecer\\ ES2 mantener};
\node[proc] (td) at (-2.6,-6.6) {\textbf{Diseño e implementación}\\ TD1--TD6};
\node[proc] (te) at (0.8,-6.6)  {\textbf{Ejecución}\\ TE1 ejecutar · TE2 comparar · TE3 registrar};
\node[proc] (ir) at (4.2,-6.6)  {\textbf{Reporte de incidentes}\\ IR1 analizar · IR2 reportar};

\draw[fl] (es) -- (td);
\draw[fl] (td) -- (te);
\draw[fl] (te) -- node[et,pos=0.5]{no aprobado} (ir);

\draw[fl] (plan.south) to[out=270,in=120] node[et,pos=0.25]{dirige} (es.north);

% ------------- bucle interno: corrección, confirmación, regresión -----
\node[salida] (nv) at (4.2,-8.5) {Nueva versión del elemento de prueba};
\draw[fl] (ir) -- node[et,pos=0.5]{corrección} (nv);

\node[conf] (pc) at (0.3,-8.5)  {\textbf{Prueba de confirmación}\\ ¿se corrigió \emph{ese} defecto?};
\node[conf] (pr) at (0.3,-10.4) {\textbf{Prueba de regresión}\\ ¿el cambio dañó lo que ya funcionaba?};
\draw[fl] (nv.west) -- (pc.east);
\draw[fl] (nv.south) to[out=250,in=340] (pr.east);

\draw[fl] (pc.north) -- node[et,pos=0.5]{nueva ejecución} (te.south);
\draw[fl] (pr.west) to[out=180,in=250,looseness=1.2] (te.south west);

% ======================= BUCLES DE RETROALIMENTACIÓN ==================
% 1) bucle de control corto: dinámica -> monitoreo -> dinámica
\draw[fb] (te.north) to[out=100,in=280] node[et,pos=0.5]{medidas y estado} (tmc.south);
\draw[fb] (tmc.west) to[out=200,in=80] node[et,pos=0.55]{directivas de control} (td.north);

% 2) bucle de replanificación: monitoreo -> planificación
\draw[fb] (tmc.north west) to[out=150,in=20] node[et,pos=0.5]{replanificación} (tp.east);

% 3) bucle de mejora organizacional: finalización -> capa organizacional
% se rodea por encima de la capa organizacional para no atravesar ningún nodo
\draw[mej] (tc.north) |- ($(ot.north)+(0,1.25)$)
             node[et,pos=0.72]{lecciones aprendidas} -- (ot.north);

% cierre hacia finalización
\draw[fl] (ir.east) to[out=20,in=250] node[et,pos=0.5]{criterios de salida} (tc.south);

% ============================ BANDAS DE CAPA ==========================
\begin{scope}[on background layer]
  \node[fit=(ot)(pol), draw=gray!45, dashed, inner sep=7pt, rounded corners] (b1) {};
  \node[fit=(tp)(plan)(tmc)(tc), draw=gray!45, dashed, inner sep=7pt, rounded corners] (b2) {};
  \node[fit=(es)(td)(te)(ir)(nv)(pc)(pr), draw=gray!45, dashed, inner sep=7pt, rounded corners] (b3) {};
\end{scope}
\node[et, text=gray!60, anchor=west, font=\tiny\bfseries] at (b1.north west) {CAPA ORGANIZACIONAL};
\node[et, text=gray!60, anchor=west, font=\tiny\bfseries] at (b2.north west) {CAPA DE GESTIÓN DE LA PRUEBA};
\node[et, text=gray!60, anchor=west, font=\tiny\bfseries] at (b3.north west) {CAPA DINÁMICA (iterativa)};

% ============================== LEYENDA ===============================
\node[align=left, font=\scriptsize, text width=14.8cm] at (0,-12.4)
  {\textbf{Cómo leer la figura.} Las flechas azules continuas indican el flujo principal; las
   naranjas discontinuas, los bucles de control y de replanificación; la verde de trazo mixto,
   el bucle lento de mejora organizacional. El monitoreo y control \emph{no} ocupa un lugar fijo
   en la secuencia: acompaña a la ejecución de principio a fin. Los pasos de la capa dinámica se
   repiten por conjunto de pruebas, incremento o iteración hasta cumplir los criterios de salida.};

\end{tikzpicture}
\caption[Flujo del proceso de prueba con sus bucles de retroalimentación]%
{Flujo del proceso de prueba en las tres capas de \normap{2}, con el monitoreo y control
representado como proceso concurrente y con los bucles de control, replanificación y mejora.}
\label{fig:procesos}
\end{figure}

\begin{observacion}[Diagrama del flujo del proceso]
\obsproblema{El diagrama de referencia rotulaba «verificar correcciones» como prueba de
regresión, se leía como una secuencia lineal, situaba el entorno y los datos dentro de la
preparación de pruebas y contemplaba sólo el análisis de riesgos del proyecto.}
\obscambio{Confirmación y regresión aparecen como cajas distintas; el monitoreo y control se
dibuja como proceso concurrente con tres bucles de retroalimentación explícitos; entorno y datos
constituyen un proceso propio (ES1/ES2); la planificación abarca riesgos de producto y de
proyecto.}
\obsfuente{OBS §5.5 y §5.6 · PDF §2.2 · DIAG}
\end{observacion}


La Figura~\ref{fig:procesos} resume el flujo del proceso de prueba en las tres capas que define
\normap{2} —organizacional, de gestión y dinámica— y, sobre todo, los bucles de
retroalimentación que impiden leerlo como una cascada: el seguimiento y control opera de forma
\emph{concurrente} con la ejecución, los incidentes y los riesgos nuevos pueden reabrir la
planificación, y las lecciones aprendidas del cierre retroalimentan la capa organizacional. Este
plan corresponde a la capa de gestión; los anexos documentan lo que ocurre en la capa dinámica.
\fi

\begin{notanormativa}[\segunvariante{Declaración de alcance de esta plantilla}{Declaración de alcance de este plan de prueba}]
\segunvariante{Esta plantilla}{Este plan} \textbf{no declara conformidad} con \norma. Organiza
la planificación y el cierre de un esfuerzo de prueba tomando como referencia el modelo de
procesos de la parte~2 y la estructura documental de la parte~3. Demostrar conformidad exigiría,
además, evidencia de que las actividades se realizaron y de que sus resultados fueron los
especificados, y una revisión del texto normativo completo de la edición elegida.
\end{notanormativa}

% ---------------------------------------------------------------------
\subsection{Referentes utilizados}

\begin{instruccion}
Relacione los documentos que sirvieron de base al plan: referencias normativas y documentación
técnica del proyecto a probar. Incluya como mínimo nombre, descripción y, para documentos
digitales, la URL y el nombre del archivo con su extensión.

Cite siempre la \emph{edición} (año) de cada norma y la \emph{versión} de cada fuente técnica.
Una referencia como «\normap{4}, Anexo~A, pp.~32--38» es ambigua sin el año, porque la
paginación cambia entre ediciones. Cuando cite una cláusula o anexo, verifíquelo en la edición
que realmente consultó.
\end{instruccion}

\begin{observacion}[Ediciones normativas]
\obsproblema{La plantilla anterior remitía a «\normap{4}, Anexo~A, páginas 32--38» sin indicar
la edición, y no pedía versionar las fuentes técnicas del proyecto.}
\obscambio{El cuadro de referentes exige edición o versión y asigna una clave que se reutiliza
para identificar la base de prueba.}
\obsjustificacion{La paginación y la numeración de anexos cambian entre ediciones, de modo que
la referencia original no es verificable.}
\obsfuente{OBS §2 y §5.2}
\end{observacion}

\begin{table}[H]
\centering
\caption{Referentes normativos y técnicos utilizados en la elaboración del plan.}
\label{tab:referentes}
\begin{tabularx}{\textwidth}{|C{1.3cm}|L{4.6cm}|Y|C{2.1cm}|}
  \hline
  \filaenc \textbf{Clave} & \textbf{Documento (con edición/versión)} & \textbf{Descripción y uso en el plan} & \textbf{Ubicación} \\ \hline
  \campo{REF-01} & \campo{Norma, edición y cláusula verificada} & \campo{Para qué se usó} & \campo{URL/archivo} \\ \hline
  \campo{REF-02} & \campo{Especificación de requisitos v1.2} & \campo{Base de prueba} & \campo{URL/archivo} \\ \hline
  & & & \\ \hline
  & & & \\ \hline
\end{tabularx}
\notatabla{Distinga los referentes \textit{normativos} (normas y estándares) de los \textit{técnicos} (documentación del proyecto a probar). Ambos son citables desde los apartados posteriores mediante su clave.}
\end{table}

El Cuadro~\ref{tab:referentes} concentra las fuentes en que se apoya el plan y les asigna una
clave; esa clave se reutiliza en el apartado~\ref{sec:objeto-base} para identificar cuáles de
ellas constituyen la \emph{base de prueba}.

% ---------------------------------------------------------------------
\subsection{Glosario}

\begin{instruccion}
Defina los términos, acrónimos y abreviaturas relevantes cuyo significado deba explicitarse.
Incluya los términos del dominio de la aplicación que faciliten comprender el caso de estudio.
Se recomienda separar acrónimos y términos en dos subapartados, en formato de lista tipo
diccionario. Indique la fuente de las definiciones tomadas de una norma o de un cuerpo de
conocimiento.
\end{instruccion}

\subsubsection*{Acrónimos}

\begin{description}[style=nextline,leftmargin=2.6cm,labelwidth=2.4cm,font=\bfseries]
  \item[\campo{SIL}] \campo{Software Integrity Level. Nivel de integridad del software; véase el apartado~\ref{sec:criticidad}.}
  \item[\campo{ACR-02}] \campo{Significado}
\end{description}

\subsubsection*{Términos}

\begin{instruccion}
Los términos siguientes se ofrecen ya redactados porque son los que con más frecuencia se
confunden entre sí al documentar pruebas. Consérvelos, ajústelos si su proyecto usa otra
terminología y añada los términos propios del dominio de la aplicación.
\end{instruccion}

\begin{observacion}[Términos predefinidos del glosario]
\obsproblema{La plantilla anterior pedía un glosario sin indicar qué términos convenía definir,
y a lo largo del documento empleaba «elemento de prueba», «cobertura» y «regresión» con
significados que se solapaban.}
\obscambio{Se predefinen los ocho términos cuya confusión originó correcciones en el resto del
documento.}
\obsfuente{OBS §4.3, §4.4 y §5.5}
\end{observacion}

\begin{description}[style=nextline,leftmargin=2.6cm,labelwidth=2.4cm,font=\bfseries]
  \item[Objeto de prueba] Elemento que se somete a prueba: sistema, subsistema, módulo,
        componente, servicio, API o interfaz. Responde a \emph{qué se prueba}.
  \item[Base de prueba] Conjunto de documentos y fuentes de los que se deriva lo que debe
        verificarse: requisitos, historias de usuario, casos de uso, reglas de negocio,
        arquitectura o diseño. Responde a \emph{contra qué se juzga el objeto}.
  \item[Condición de prueba] Aspecto verificable del objeto de prueba, derivado de la base de
        prueba y de los riesgos, a partir del cual se diseñan casos.
  \item[Cobertura] Proporción de elementos de un conjunto identificado —requisitos,
        particiones, decisiones, transiciones— que han quedado \emph{ejercitados} por las
        pruebas. Se mide con independencia del veredicto obtenido: un elemento probado por un
        caso que falló está cubierto. No debe confundirse con el \emph{avance de ejecución}
        (cuánto del trabajo previsto se ha realizado) ni con el \emph{veredicto} (si lo
        ejercitado se comportó como se esperaba). Véase el apartado~\ref{sec:cobertura}.
  \item[Veredicto] Juicio emitido al comparar el resultado real de una ejecución con el
        esperado: aprobado, no aprobado o no concluyente. Pertenece a la ejecución, no al caso
        de prueba.
  \item[Defecto] Imperfección en un producto de trabajo que puede hacer que no cumpla sus
        requisitos. Es estático y reside en el artefacto.
  \item[Fallo] Evento en el que el servicio entregado se desvía del correcto; es la
        manifestación externa observable de un defecto activado.
  \item[Incidente de prueba] Evento ocurrido durante la prueba que requiere investigación o
        acción adicional. Es un registro documental: un incidente puede deberse a un defecto
        del objeto, del entorno, de los datos o del propio caso de prueba.
  \item[\campo{Término del dominio}] \campo{Definición}
\end{description}

\begin{notanormativa}[Fuente de las definiciones anteriores]
Las definiciones de defecto, fallo e incidente siguen la delimitación empleada en
ISO/IEC/IEEE 24765 y en \normap{1}, recogida en la sección~3.1 del documento
\emph{Plan de trabajo 29119} incluido en \texttt{sources/}. Un mismo defecto puede producir
varios fallos distintos y un mismo fallo puede tener varios defectos como causa.
\end{notanormativa}
